\section{超越二人零和：Population Best Response、Ultimatum 与人类数据必要性}

\subsection{核心转折：多方博弈下 Minimax 失去解释力}
进入非二人零和后，课程结论很明确：``求一个 Minimax'' 不再是有意义目标。关键问题转为 ``你要对哪个群体最优''。这个群体可能是同类 Agent，也可能是人类玩家，二者分布差异巨大。

\begin{importantbox}{Brown 的强主张}
若目标是与人类合作并稳定获益，避免使用 human data 基本是死路（dead end）。
\end{importantbox}

\subsection{Ultimatum Game：同一规则，不同文化分布}
在 Ultimatum Game 中，理论 Nash 均衡与人类真实行为常显著偏离。课堂中给出的经验范围是：很多人会在 20\%--30\% 以上才愿意接受报价，过低报价即使理性上 ``应接受'' 也会被拒绝。

\begin{knowledgebox}{这对 Agent 训练的含义}
没有目标人群行为数据，模型只能学到 ``在自博弈里自洽''，却学不到 ``在真实人群里可协作''。前者可达成，后者不可自动涌现。
\end{knowledgebox}

\subsection{Diplomacy 案例：DORA 与 SearchBot 的交叉失配}
课程中最有说服力的证据来自 Diplomacy：
\begin{itemize}
    \item DORA：纯 self-play 训练，在某些设定里可达超人；
    \item SearchBot：人类数据驱动策略，在自身群体里表现更稳。
\end{itemize}
当两类策略群体交叉对战时，出现明显分布错配：各自都可能在 ``非本群体'' 环境中退化。这证明了多均衡共存与群体依赖性。

\begin{warningbox}{策略评估的常见错误}
只在 ``同类自博弈池'' 里评估就宣布 SOTA，往往会高估跨群体泛化。对外部署前必须做跨群体、跨文化、跨协议评测。
\end{warningbox}

\subsection{可操作训练配方}
Brown 给出的配方可总结为三步：
\begin{enumerate}
    \item 收集目标群体数据并训练 imitation model；
    \item 放大 inference-time compute，以更准确建模群体行为；
    \item 让 RL 在 ``人类仿真群体环境'' 中继续优化。
\end{enumerate}

这是 Cicero 及相关系统成功路径的抽象版本。其本质不是放弃 RL，而是让 RL 在正确分布上优化。

\subsection{本章小结}
非二人零和任务里，``先定义群体，再定义最优'' 是必要顺序。没有群体数据，Population Best Response 无法落地。

